%================================================
% CAPÍTULO ARRAYS
%================================================

\chapterimage{chapter_pointers.jpg} % Chapter heading image

\chapter{Arrays}
Quem nunca precisou organizar dados de tipo semelhante num único local? Eu sim, você não? Então quando você, por exemplo, escreve um texto num pedaço de papel, o que acha que está fazendo? \say{\textit{Organizando ideias?}}. Sim, também! Contudo, o que você está organizando materialmente são caracteres.

Imagine agora que, usando o mesmo texto do papel no computador, para cada caractere organizado tenhamos que definir uma variável do tipo \textit{char} e atribuir-lhe o valor respectivo.

\begin{ccode}
  char c1, c2, c3, c4, c5, c6, c7, c8, c9, c10,
       c11, c12, c13, c14, c15, c16, c17;

  c1 = 'A';
  c1 = 'u';
  c1 = 't';
  c1 = 'o';
  c1 = 'r';
  c1 = ' ';
  ...
\end{ccode}
\\

Tal exemplo é deveras trabalho do capiroto, e só será utilizado na vida real se o programador tiver algum pacto com esse tinhoso. Para os que seguem o caminho da luz e dos ensinamentos jedi há uma alternativa feliz.

\textbf{Arrays}, em C, são tipos de dados de alocação contíngua derivados dos tipos simples. Em outras palavras, são estruturas de dados homogêneas, compostas de elementos de mesmo tipo, e posicionados em sequência na memória. Arrays também são conhecidos como \textit{vetores}, quando a sequência dos dados é unidimensional (linear), e como \textit{matrizes}, quando essa sequência se dá multimensionalmente (linhas e colunas). Divagaremos a respeito dos arrays do tipo matriz mais a frente, abordando, por enquanto, o array do tipo vetor. Seguem dois exemplos de arrays.

\begin {center}
  \begin{bytefield}[bitwidth=1em]{17}
    \bitbox{17}{Autor presunçoso!}\\
    \bitbox{1}{A} \bitbox{1}{u} \bitbox{1}{t} \bitbox{1}{o} \bitbox{1}{r} \bitbox{1}{ }
    \bitbox{1}{p} \bitbox{1}{r} \bitbox{1}{e} \bitbox{1}{s} \bitbox{1}{u} \bitbox{1}{n}
    \bitbox{1}{ç} \bitbox{1}{o} \bitbox{1}{s} \bitbox{1}{o} \bitbox{1}{!}
  \end{bytefield}
\end{center}

Nesse primeiro, a frase \say{Autor presunçoso!} é a organização de dezessete caracteres um atrás do outro. Isso mesmo, dezessete, pois o espaço também é um dado. Nunca se esqueça que, quando tratamos dados, até um \say{0} ocupa espaço na memória do computador.

Abaixo outra organização de dados. Desta vez de inteiros que serão o resultado do próximo sorteio da mega-sena. Duvida?

\begin{center}
  \begin{bytefield}[endianness=little,bitwidth=6em]{6}
    \bitbox{1}{9} \bitbox{1}{20} \bitbox{1}{11}
    \bitbox{1}{7} \bitbox{1}{43} \bitbox{1}{31}\\
    % \bitheader{0-5}
  \end{bytefield}
\end{center}

\begin{remark}
Em algumas fontes esparsas em língua portuguesa, você, leitor, encontrará referências a arrays como sendo \textit{\textbf{arranjos}}. Ressaltamos que um array, numa tradução mais fiel, é um \textit{\textbf{conjunto}} de dados. Já um arranjo:
  \begin{figure}[!htp]
    \centering
    \includegraphics[scale=0.8, keepaspectratio=true]{arrays_arrange.jpg}
    \caption{Array não é Arranjo}
    \label{fig:arrays_arrange}
  \end{figure}
\end{remark}

\section{Usando Arrays}
\subsection{Definindo Arrays}

Ao definirmos um array, precisamos determinar a respectiva quantidade de elementos.

\begin{ccode}
  int i[6];
  char c[3];
\end{ccode}
\\

Os colchetes, \textit{operadores} de índice \textbf{[]}, informam ao compilador que as variáveis \textbf{i} e \textbf{c} são do tipo array. Já o número presente entre esses operadores define a quantidade de elementos a serem alocados para o array em questão. No código acima, \textbf{i} é definido como um array de \textit{int} com seis elementos e \textbf{c} como um array de \textit{char} com três elementos.

\begin{center}
  \begin{bytefield}[endianness=little,bitwidth=6em]{6}
    i[6]\\
    \bitbox{1}{-57} \bitbox{1}{3} \bitbox{1}{791431480}
    \bitbox{1}{0} \bitbox{1}{32767} \bitbox{1}{0}\\
    \bitheader{0-5}\\
  \end{bytefield}
  \begin{bytefield}[endianness=little,bitwidth=2em]{3}
    c[3]\\
    \bitbox{1}{0} \bitbox{1}{-10} \bitbox{1}{49}\\
    \bitheader{0-2}
  \end{bytefield}
\end{center}

Uma boa prática de programação é inicializar os arrays, assim como os tipos simples, pois, tanto em i[] quanto em c[], vemos claramente que os elementos têm valores aleatórios. Isso se dá pelo lixo encontrado na área da memória usada por esses recém-definidos arrays num escopo local.

Arrays também podem ser definidos com a quantidade de elementos informada por outra variável, quando em escopo local. Mas atente sempre para inicializar a variável que é responsável pela quantidade de elementos acima de zero, do contrário comportamentos indesejados acontecerão, tais como pipocas voando na tela ou suco de tangerina escorrendo pelos cantos do monitor.

\textit{Inicialize sempre a variável com a quantidade de elementos.}\\
\begin{ccode}
  int elements = 10;
  ...
  int i[elements];
\end{ccode}
\\

\textit{Nunca se esqueça de inicializar! NUNCA! Ou o bicho vem pegar você...}\\

\subsection{Declarando Arrays}
Faz bem relembrarmos que declarar não é o mesmo que definir, portanto, ao contrário do que se pensa comumente, arrays podem normalmente ser declarados. Para tanto, basta fazermos uso do modificador \textit{extern} já visto no item 3.3.

Forçando a quantidade de elementos \textbf{ext1[32]} ou não \textbf{ext2[]}.

\textit{array.h - Declaração}\\
\begin{ccode}
extern int ext1[32];
extern int ext2[];
\end{ccode}
\\

E posteriormente definindo os mesmos arrays

\textit{main.c - Definição}\\
\begin{ccode}
#include "array.h"
...
int ext1[32];
int ext2[32];
...
\end{ccode}


\subsection{Inicializando Arrays}
Recordar é viver, então viva mais uma vez: \say{inicializar é atribuir valor no ato da declaração}. No caso de arrays podemos inicializá-los de várias formas, vejamos.
\\

Inicializando os elementos um a um.

\begin{ccode}
  int arr[3] = {1, 5, 10};
\end{ccode}
\\

Inicializando os elementos um a um e forçando o compilador a reservar a quantidade de elementos informados na inicialização. Neste caso, o array é declarado e inicializado com 4 elementos.

\begin{ccode}
  int arr[] = {3, 2, 9, 1};
\end{ccode}
\\

Inicializando todos os elementos com zero.

\begin{ccode}
  int arr[3] = {0};
\end{ccode}
\\

Inicializando todos os elementos com zero (modo \say{preguiçoso suvina}, economiza um byte no código fonte por inicialização).

\begin{ccode}
  int arr[3] = {};
\end{ccode}
\\

Inicializando o primeiro elemento com 5 e os demais com zero.

\begin{ccode}
  int arr[3] = {5};
\end{ccode}
\\

Quando um array é definido num escopo global \textit{e/ou} com a classe de armazenamento \textit{static} e não há inicialização explícita, de acordo com o \textit{C Standard}, todos os elementos são inicializados com zero.

Inicializando todos os elementos com zero por definiçao estática.

\begin{ccode}
  static int array[6];
\end{ccode}
\\

Inicializando por escopo global.

\begin{ccode}
  int arr[6];  /* Os elementos de arr iniciam
                  com zero, automaticamente.
                  "static" subentendido */

  int main(void)
  {
    return 0;
  }
\end{ccode}
\\

Há, também, um array especial chamado array flexível. Ele só pode ser usado como último membro de uma estrutura que já contém pelo menos outro membro. Esse array se destina ao uso em estruturas com tamanho dinâmico que precisam ser compartilhadas integralmente e/ou acessadas linearmente.

Veja abaixo exemplo de uso de um array flexível.

\begin{ccode}
  struct str {
      int body_len;
      double body[];
  };

  int length = sizeof(double [10]);

  struct str *s = malloc(sizeof(struct str) + length);
  s->body_len = length;
  ...
\end{ccode}
\\

E há o array \say{sou brasileiro, não desisto nunca} que, tal qual o mote, não tem sentido na vida real. Esse array aloca apenas um elemento.

\begin{ccode}
  int brazilian_i[] = {1};
  float brazilian_f[1] = {2};
\end{ccode}
\\

\subsection{Acessando Arrays}
Os elementos dos arrays em C são acessados através do índice 0 até a quantidade de elementos menos um (n-1), ou seja, num array com três elementos podemos acessá-los da seguinte maneira.

\begin{ccode}
  float att[3] = {10.5, 25.3, 9.0};

  printf("%f\n", att[0]);
  printf("%f\n", att[1]);
  printf("%f\n", att[2]);
\end{ccode}
\\

O código acima imprime na tela os valores dos elementos att[0] que é 10.5, att[1] com 25.3, e o att[2] respectivamente 9.0. O array como qualquer variável se encontra em locais específicos da memória. Antes do primeiro e logo após seu último elemento há dados referentes a outros controles ou estruturas do próprio programa, por isso, acessá-lo, inadvertidamente, além da área definida na declaração/inicialização ocasionará seguramente comportamentos inesperados no software (bugs), até porque o compilador não o avisará de plano sobre limites extrapolados se você não estiver utilizando flags de warning (-Warray-bounds, neste caso). Portanto, atenção redobrada e vergonha na cara.

\subsubsection{Atribuindo Valores em Arrays}
Até então vimos apenas como informar ao compilador o quanto de espaço precisamos que ele reserve na memória para nossos dados e como atribuir valores através da inicialização. \textit{Vamos bagunçar o coreto, então?}

\say{Pensei que não iria falar nada sobre atribuição...}

Foi mesmo, bateu uma preguiça... /o{\textbackslash} A verdade é que discorrer brevemente o tema índices no tópico anterior, também conhecido como subscritores, foi crucial para mostrar a você, aspirante a programador C, que nesta linguagem, assim como grande parte das demais, o índice inicial de conjunto de dados é o 0 e não o 1. Sem esse conhecimento básico, talvez você já tivesse desistido de C e iniciado a programar em Pascal.

\say{Mais falatório e não explicou como atribuo valores num array!}

Nervosinho, você, hein? Calma que é simples.

\begin{ccode}
  float att[3];

  att[0] = 1.1;
  att[1] = 1.2;
  att[2] = 1.3;
\end{ccode}
\\

Pronto! Esse último código não inicializa o array deixando essa função para a atribuição singular de cada elemento. Doeu tanto esperar?

\section{Array de Caracteres}

\section{Arrays na Vida Real?!}
\subsection{Imprimir Elementos}
\subsection{Somar Elementos}
\subsection{Inverter o Array}
\subsection{Ordenar o Array}

\section{Arrays como Parâmetros}

\section{Arrays Multidimensionais (Matrizes)}


%------------------------------------------------
